
{\actuality} В настоящее время актуальной является проблема группового применения роботизированных наземных транспортно-технологических машин (РНТТМ), функционирующих автономно и эффективно выполняющих задачи по предназначению. Решение указанной проблемы требует разработки методов управления, таких как многозадачную координацию в группе машин. Указанные методы для эффективного функционирования РНТТМ перспективно строить в рамках децентрализованного управления, с возможной минимизацией объема информации, которой машины обмениваются. Это обусловлено рядом ограничений при применении метода централизованного управления РНТТМ: во-первых, повышенные риски к потере управления группой машин при выходе из строя центрального управляющего лидера \cite{habr2012kvis} или прерываниях в каналах связи, а~во-вторых, время принятия решения группой машин экспоненциально увеличивается с~увеличением числа объектов в группе. В~децентрализованной системе группового управления РНТТМ нет четкого лидера, а каждый член группы РНТТМ автономен в той или иной мере.

Различным аспектам проблемы группового управления РНТТМ посвящены работы отечественных
(Алешин~М.\,В., Гайдук~А.\,Р., Довгаль~В.\,А., Журба~Я.\,С., Казарян~Д.\,Э., Козлова~М.\,Г., Лазарев~В.\,С., Медведев~М.\,Ю., Певзнер~Л.\,Д., Петренко~В.\,И., Пшихопов~В.\,Х., Феофилова~А.\,А., Фильченков~А.\,А., Чичерин~И.\,В. и~др.) и зарубежных (Aziz~H., Braquet~M., Choudhury~S., Dias~M.\,B., Gerkey~B.\,P., Korsah~G.\,A., Liu~W., Mailler~R., Nunes~E., Rizk~Y., Sawadsitang~S., Velhal~S., Wei~Y., Yan~Z., Zhang~F. и~др.) ученых.

Из анализа этих работ можно сделать вывод, что наиболее перспективными являются групповые системы с децентрализованным типом управления РНТТМ. Такой подход предполагает взаимодействие нескольких машин и является децентрализованным, поскольку в нем предполагается коллективное управление системой машин на основе решений, вырабатываемых всеми членами группы машин. Однако такие системы управления требуют интенсивного обмена информацией между машинами, что является проблемой для подсистемы связи, которая должна обеспечивать высокую помехоустойчивость и оперативную реконфигурацию системы машин.

% {\progress}
% Этот раздел должен быть отдельным структурным элементом по
% ГОСТ, но он, как правило, включается в описание актуальности
% темы. Нужен он отдельным структурынм элемементом или нет ---
% смотрите другие диссертации вашего совета, скорее всего не нужен.

{\aim} данной работы является разработка метода автоматизации по многозадачной координации между роботизированными наземными транспортно-технологическими машинами (РНТТМ), направленного на повышение их производительности и энергоэффективности.

\textbf{Актуальность работы} обусловлена необходимостью разработки теоретических и методологических средств управления эксплуатацией РНТТМ, направленных на повышение качества выполнения дорожно-строительных работ с соблюдением технологических стандартов и, как следствие, достижение целевых показателей безопасности и состояния дорожной сети, установленных Указом Президента Российской Федерации от 7 мая 2024 года № 309 <<О национальных целях развития Российской Федерации на период до 2030 года и на перспективу до 2036 год>> \cite{actPRF2024}. Актуальность также определена наличием проблемы эффективного группового применения РНТТМ, функционирующих автономно и эффективно выполняющих задачи по предназначению.

Так, в настоящее время стратегические ориентиры развития дорожной инфраструктуры и обеспечения транспортной безопасности определяются новыми национальными целями, установленными Указом ПРФ \cite{actPRF2024}. В соответствии с этим документом, одной из приоритетных задач в рамках национальной цели <<Комфортная и безопасная среда для жизни>> является:

\begin{itemize}
	\item увеличение доли соответствующих нормативным требованиям автомобильных дорог федерального значения и дорог крупнейших городских агломераций не менее чем до 85\% к 2030 году (пункт 4(л) \cite{actPRF2024});

	\item снижение смертности в результате дорожно-транспортных происшествий в полтора раза к 2030 году и в два раза к 2036 году (пункт 4(м) \cite{actPRF2024}).
\end{itemize}

Достижение указанных показателей требует принципиального обновления методологии организации и управления процессами эксплуатации НТТМ при строительстве, реконструкции и содержании автомобильных дорог общего пользования. Ключевую роль в этом играет научное обоснование системы управления процессом обеспечения эффективной эксплуатации наземных транспортно-технологических машин, используемых в дорожной отрасли.

Для~достижения поставленной цели необходимо решить следующие {\tasks}:
\begin{enumerate}[beginpenalty=10000] % https://tex.stackexchange.com/a/476052/104425
  \item Провести анализ известных методов автоматизации рабочих процессов РНТТМ.
  \item Определить, формализовать и разработать имитационную модель многозадачной координации в~группе РНТТМ.
  \item Разработать алгоритмы многозадачной координации в группе РНТТМ на основе модели.
  \item Разработать метод автоматизированной многозадачной координации в~группе РНТТМ.
\end{enumerate}

\textbf{Предметом исследования} является группа роботизированных наземных транспортно-технологических машин.

\textbf{Объектом исследования} является их многозадачная координация при выполнении задач по предназначению эксплуатации группы машин.

{\novelty}
\begin{enumerate}[beginpenalty=10000]
  \item Проведен анализ известных методов автоматизации рабочих процессов РНТТМ.
  \item Разработана имитационная модель многозадачной координации для группы роботизированных наземных транспортно-технологических машин, в которой учтены особенности энергоограниченности и неоднородности оборудования, что позволяет повысить производительность и энергоэффективность.
  \item Разработаны алгоритмы стратегий координации (энергосберегающая стратегия и минимизации временных задержек), основанные на адаптации методов теории игр (включая байесовские равновесия) к задачам распределения ресурсов в условиях стохастической постановки задач.
  \item Разработан метод автоматизированной многозадачной координации в~группе РНТТМ и программная платформа для моделирования и сравнительного анализа различных стратегий координации в условиях, приближенных к реальной эксплуатации (учет стохастичности задач, энергопотребления и типа оборудования).
\end{enumerate}

{\influence} заключается в применении разработанного метода многозадачной координации в группе РНТТМ !!! -> из 4й главы формулировки вынести сюда

{\methods} При выполнении работы были использованы методы системного анализа, математического моделирования и~тестирования программного обеспечения, теории вероятностей, игр и~программирования высокого уровня, статистические, а так же численные методы.

{\defpositions}
\begin{enumerate}[beginpenalty=10000]
	\item Имитационная модель многозадачной координации группы роботизированных НТТМ с учетом энергоограничений и~стохастической природы задач.
	\item Алгоритмы и стратегии координации задач в группе РНТТМ, обеспечивающие повышение эффективности распределения задач в условиях дефицита ресурсов.
	\item Метод автоматизированной многозадачной координации в~группе РНТТМ, основанный на энергосберегающей стратегии, который обеспечивает большее количество успешно выполненных задач по сравнению с другими методами, что подтверждает ее практическую применимость.
\end{enumerate}

\textbf{Область диссертационного исследования} -- соответствует паспорту научной специальности 2.5.11. Наземные транспортно-технологические средства и комплексы, пункту 3 -- <<Экспериментальные исследования и испытания транспортно-технологических средств и их комплексов, а также отдельных систем, агрегатов, узлов, деталей и технологического оборудования>>, пункту 5 -- <<Математическое моделирование рабочих процессов транспортно-технологических средств, в том числе в их узлах, механизмах, системах и технологическом оборудовании при взаимодействии с опорной поверхностью и с рабочими средами (объектами)>>.

{\reliability} полученных результатов обеспечивается проведением серии вычислительных экспериментов с многократными прогонами разработанного метода автоматизированной многозадачной координации в различных сценариях, моделирующих реальные условия эксплуатации машин. Статистическая значимость результатов подтверждается применением методов математической статистики для обработки экспериментальных данных, включая вычисление выборочных средних значений и стандартных отклонений ключевых метрик производительности. Полученные результаты подтверждают эффективность предложенного метода, показывая устойчивое превосходство энергосберегающей стратегии по ключевым показателям (количество выполненных задач, энергоэффективность, время выполнения) по сравнению с аналогами, что находится в соответствии с результатами, полученными в ходе вычислительных экспериментов.

{\probation}
Основные результаты работы докладывались~на:

\begin{enumerate}[beginpenalty=10000]
	\item 28-й Московской международной межвузовской научно-технической конференции <<Подъемно-транспортные, строительные, дорожные, путевые, мелиоративные машины и робототехнические комплексы>> (<<Молодой инженер>>) (МГТУ им.~С.М.Баумана, Москва, 2024\,г.),
	\item XVI Международной конференции <<Организация и безопасность дорожного движения в крупных городах>>, посвященной памяти доктора технических наук, профессора Кравченко Павла Александровича (СПб ГАСУ, Санкт-Петербург, 2024\,г.),
	\item XXXVIII Национальной (с международным участием) научно-технической конференции <<Эксплуатация и сервис автомобилей, тракторов и двигателей>> (СПб ГАУ, Санкт-Петербург, 2025\,г.),
	\item 29-й Московской международной межвузовской научно-технической конференции <<Подъемно-транспортные, строительные, дорожные, путевые, мелиоративные машины и~робототехнические комплексы>> (<<Молодой инженер>>) (МГТУ им.~С.М.Баумана, Москва, 2025\,г.),
	\item VII Международной научно-практической конференции <<Мехатроника, автоматика и~ робототехника>> (НИЦ МС, Москва, 2025\,г.).
\end{enumerate}

Результаты работы докладывались на научных семинарах кафедры <<Наземных транспортно-технологических средств и комплексов>> Санкт-Петербургского государственного архитектурно-строительного университета (2022-2026 гг.).

{\contribution} Автором лично:
\begin{enumerate}[beginpenalty=10000]
    \item Предложена функция полезности для учета энергетических ограничений при выборе задач;
    \item Формализованы и реализованы алгоритмы энергосберегающей стратегии и минимизации временных задержек;
    \item Разработана и реализована программная платформа моделирования и~анализа стратегий координации;
    \item Проведена серия вычислительных экспериментов и статистический анализ результатов;
    \item Результаты работы внедрены в войсковых частях 13821 и 31895, а также в учебном процессе кафедры НТТСиК СПб ГАСУ.
\end{enumerate}

\ifnumequal{\value{bibliosel}}{0}
{%%% Встроенная реализация с загрузкой файла через движок bibtex8. (При желании, внутри можно использовать обычные ссылки, наподобие `\cite{vakbib1,vakbib2}`).
    {\publications} Основные результаты по теме диссертации изложены
    в~XX~печатных изданиях,
    X из которых изданы в журналах, рекомендованных ВАК,
    X "--- в тезисах докладов.
}%
{%%% Реализация пакетом biblatex через движок biber
    \begin{refsection}[bl-author, bl-registered]
        % Это refsection=1.
        % Процитированные здесь работы:
        %  * подсчитываются, для автоматического составления фразы "Основные результаты ..."
        %  * попадают в авторскую библиографию, при usefootcite==0 и стиле `\insertbiblioauthor` или `\insertbiblioauthorgrouped`
        %  * нумеруются там в зависимости от порядка команд `\printbibliography` в этом разделе.
        %  * при использовании `\insertbiblioauthorgrouped`, порядок команд `\printbibliography` в нём должен быть тем же (см. biblio/biblatex.tex)
        %
        % Невидимый библиографический список для подсчёта количества публикаций:
        \phantom{\printbibliography[heading=nobibheading, section=1, env=countauthorvak,          keyword=biblioauthorvak]%
        \printbibliography[heading=nobibheading, section=1, env=countauthorwos,          keyword=biblioauthorwos]%
        \printbibliography[heading=nobibheading, section=1, env=countauthorscopus,       keyword=biblioauthorscopus]%
        \printbibliography[heading=nobibheading, section=1, env=countauthorconf,         keyword=biblioauthorconf]%
        \printbibliography[heading=nobibheading, section=1, env=countauthorother,        keyword=biblioauthorother]%
        \printbibliography[heading=nobibheading, section=1, env=countregistered,         keyword=biblioregistered]%
        \printbibliography[heading=nobibheading, section=1, env=countauthorpatent,       keyword=biblioauthorpatent]%
        \printbibliography[heading=nobibheading, section=1, env=countauthorprogram,      keyword=biblioauthorprogram]%
        \printbibliography[heading=nobibheading, section=1, env=countauthor,             keyword=biblioauthor]%
        \printbibliography[heading=nobibheading, section=1, env=countauthorvakscopuswos, filter=vakscopuswos]%
        \printbibliography[heading=nobibheading, section=1, env=countauthorscopuswos,    filter=scopuswos]}%
        %
        \nocite{*}%
        %
        {\publications} Основные результаты по теме диссертации изложены в~\arabic{citeauthor}~печатных изданиях,
        \arabic{citeauthorvak} из которых изданы в журналах, рекомендованных ВАК%
        \ifnum \value{citeauthorscopuswos}>0%
            , \arabic{citeauthorscopuswos} "--- в~периодических научных журналах, индексируемых Web of~Science и Scopus%
        \fi%
        \ifnum \value{citeauthorconf}>0%
            , \arabic{citeauthorconf} "--- в~тезисах докладов.
        \else%
            .
        \fi%
        \ifnum \value{citeregistered}=1%
            \ifnum \value{citeauthorpatent}=1%
                Зарегистрирован \arabic{citeauthorpatent} патент.
            \fi%
            \ifnum \value{citeauthorprogram}=1%
                Зарегистрирована \arabic{citeauthorprogram} программа для ЭВМ.
            \fi%
        \fi%
        % Сколько патентов и программ
        \ifnum \value{citeregistered}>1%
            Зарегистрированы\ %
            \ifnum \value{citeauthorpatent}>0%
            \formbytotal{citeauthorpatent}{патент}{}{а}{}%
            \ifnum \value{citeauthorprogram}=0 . \else \ и~\fi%
            \fi%
            \ifnum \value{citeauthorprogram}>0%
            \formbytotal{citeauthorprogram}{программ}{а}{ы}{} для ЭВМ.
            \fi%
        \fi%
        % К публикациям, в которых излагаются основные научные результаты диссертации на соискание учёной
        % степени, в рецензируемых изданиях приравниваются патенты на изобретения, патенты (свидетельства) на
        % полезную модель, патенты на промышленный образец, патенты на селекционные достижения, свидетельства
        % на программу для электронных вычислительных машин, базу данных, топологию интегральных микросхем,
        % зарегистрированные в установленном порядке.(в ред. Постановления Правительства РФ от 21.04.2016 N 335)
    \end{refsection}%
    \begin{refsection}[bl-author, bl-registered]
        % Это \dfrac{refsection}{знамен}=2.
        % Процитированные здесь работы:
        %  * попадают в авторскую библиографию, при usefootcite==0 и стиле `\insertbiblioauthorimportant`.
        %  * ни на что не влияют в противном случае
        \nocite{vakbib1}%vak
        % \nocite{patbib1}%patent
        \nocite{progbib1}%program
        % \nocite{bib1}%other
        \nocite{confbib1}%conf
    \end{refsection}%
        %
        % Всё, что вне этих двух refsection, это refsection=0,
        %  * для диссертации - это нормальные ссылки, попадающие в обычную библиографию
        %  * для автореферата:
        %     * при usefootcite==0, ссылка корректно сработает только для источника из `external.bib`. Для своих работ --- напечатает "[0]" (и даже Warning не вылезет).
        %     * при usefootcite==1, ссылка сработает нормально. В авторской библиографии будут только процитированные в refsection=0 работы.
}
